\begin{helper}
Assume the parameter hierarchy.  Let
\[
  K=K(n)=\operatorname{TwoBiteNaturalIndependenceScale}(1+\epsilon,n)
\]
and let
\[
  S=\left\lceil
    \frac{K n^{1/4-\epsilon}+L_{\mathrm{large}}(\epsilon,n)}
         {L_{\mathrm{small}}(\epsilon,n)}
  \right\rceil .
\]
Eventually in \(n\),
\[
  K\log n+2S\log n
    \le K\log n+3K n^{1/4-3\epsilon}\log n .
\]
Moreover,
\[
  \exp\!\left(
    K\log n+3K n^{1/4-3\epsilon}\log n
      -n^{3/4-2\epsilon}
  \right)\longrightarrow 0 .
\]
\end{helper}
\begin{proof}
The ceiling bound from \noderef{MediumClosedPairsWitnessSizeCeilPackage} gives
\[
  S\le K n^{1/4-3\epsilon}+n^{1/4-\epsilon}+1 .
\]
By \noderef{ParameterHierarchyBaseThreshold}, eventually
\[
  K\ge (1+\epsilon)\sqrt{n\log n}.
\]
Together with \(\epsilon<1/16\), this implies, for all sufficiently large
\(n\),
\[
  n^{1/4-\epsilon}+1
    \le \frac12 K n^{1/4-3\epsilon}.
\]
Thus \(S\le \frac32 K n^{1/4-3\epsilon}\).  Since \(\log n\ge 0\)
eventually, multiplying by \(2\log n\) gives the displayed count-exponent
bound.

For the exponential envelope, use \noderef{ParameterHierarchyT28Threshold}
and \noderef{ParameterHierarchyBaseThreshold}.  The first logarithmic term is
at most \(\frac14 n^{3/4-2\epsilon}\), while the coefficient-three middle
term is at most \(\frac38 n^{3/4-2\epsilon}\).  Hence the exponent is at most
\[
  -\frac38 n^{3/4-2\epsilon}.
\]
Because \(3/4-2\epsilon>0\), this upper bound tends to \(-\infty\), and its
exponential tends to \(0\).  The original exponential is nonnegative and
bounded above by this vanishing envelope, so it also tends to \(0\).
\end{proof}
